\uuid{B5GX}
\titre{Intégration/DL en l'infini}
\auteur{hanani ; mzari}
\organisation{exo7}
\licence{CC BY-NC-SA 4.0}
\datecreate{2020-01-20}
\niveau{L1}
\module{Analyse}
\chapitre{Développement limité}
\sousChapitre{Calculs}
\difficulte{3}
\theme{Opérations sur les DL, Intégration, DL en l'infini}
\exo7qcmid{lille-2.developpements-limites.019}

\contenu{
\texte{
Soit $a\in \Rr$ et $f$ une fonction. On dit que $f(x)$ est un $o\left((x-a)^n\right)$ au voisinage de $a$ et on écrit $f(x)=o((x-a)^n)$ si $\displaystyle \lim _{x\to a}\frac{f(x)}{(x-a)^n}=0$. On note par $DL_n(a)f(x)$ le développement limité de $f$ à l'ordre $n$ en $a$. On note aussi $DL_n(a^+)f(x)$ (resp. $DL_n(a^-)f(x)$) le développement limité de $f$ à l'ordre $n$ à droite en $a$ (resp. à gauche en $a$).
}

\begin{question}
\qclassification{125.02}
\qtags{motcle=Intégration/DL en l'infini}
Soit $\displaystyle f(x)=\arctan x$. Pour $t\neq 0$, on pose $\displaystyle g(t)=f\left(\frac{1}{t}\right)$. Parmi les affirmations suivantes, cocher celles qui sont vraies :
\begin{answers}
\bad{Pour $t>0$, on a : $\displaystyle g(t)=\frac{\pi}{2}+\arctan (t)$.}
\good{Pour tout $x>0$, $\displaystyle g'(t)=\frac{-1}{1+t^2}$.}
\bad{$\displaystyle DL_3(0^+)g(t)=-t+\frac{t^3}{3}+o(t^3)$.}
\good{$\displaystyle DL_3(+\infty)f(x)=\frac{\pi}{2}-\frac{1}{x}+\frac{1}{3x^3}+o\left(\frac{1}{x^3}\right)$.}
\end{answers}

\vskip2mm

\begin{explanations}
Pour $t>0$, on a : $\displaystyle g(t)=\arctan \left(\frac{1}{t}\right)=\frac{\pi}{2}-\arctan (t)$. Donc
$$\displaystyle g'(t)=-\frac{1}{1+t^2}=-1+t^2+o(t^2)$$
et, par intégration, $\displaystyle DL_3(0^+)g(t)=\frac{\pi}{2}-t+\frac{t^3}{3}+o(t^3)$ car $\arctan 0=0$. Ainsi
$$\displaystyle DL_3(+\infty)f(x)=\frac{\pi}{2}-\frac{1}{x}+\frac{1}{3x^3}+o\left(\frac{1}{x^3}\right).$$
\end{explanations}
\end{question}
}
